\subsection{Structural necessity of non-injective projection}
  \label{subsec:structural-necessity}

  The spectral derivation developed above shows that the renormalized variation
  of the projective entropy functional produces an infrared-dominant Einstein
  response.
  It is useful to clarify under which structural conditions such an emergent
  gravitational sector can arise.

  Consider a physical description obtained through a projection
  \[
    \Pi : F \rightarrow O ,
  \]
  mapping a space of admissible configurations $F$ to a space of observable
  states $O$.
  The space $F$ is not an external substrate: it is induced by the admissibility
  structure of the observable states, with its internal organisation constrained
  by the Born\textendash Infeld capacity bound~\cite{Beau2026Foundation,Beau2026l}.
  Such projections need not be injective: distinct admissible configurations may
  yield the same observable outcome.

  A basic structural observation is that admissible injective projections possess
  trivial holonomy.
  If the projection is injective, transport along closed loops in the projection
  fiber is necessarily trivial.
  The associated fiber connection therefore has vanishing curvature,
  \[
    \Omega_{\mu\nu} = 0 .
  \]

  In this case the projection cannot generate additional geometric structure in
  the effective description.
  In particular, no contribution to the effective spacetime curvature can arise
  from projection-induced fiber curvature.

  When the projection is non-injective, by contrast, the projection fiber may
  carry non-trivial curvature.
  The effective geometric response may then contain terms schematically of the
  form
  \[
    R^{\mathrm{eff}}_{\mu\nu}
    \supset
    \mathrm{tr}_{g_\Pi}\!\left(\Omega_{\mu\alpha}\Omega_\nu{}^{\alpha}\right),
  \]
  showing that curvature in the projection fiber can act as a source for
  effective spacetime curvature.

  This observation provides a structural interpretation of emergent gravitation.
  If the effective description were injective, the projection would introduce no
  additional curvature and no gravitational sector would arise from projection
  structure.
  The existence of an emergent gravitational response therefore requires a
  non-injective projection~\cite{Beau2026l}.


  \begin{remark}[Heisenberg torsion as the structural source]
    The non-commutativity underlying the gravitational response has a concrete
    realisation in the Cosmochrony spectral programme.
    In $\mathrm{Heis}_3(\mathbb{Z}/q\mathbb{Z})$, the generating set
    $S_q = \{X^{\pm 1}, Y^{\pm 1}\}$ does not include the central generator $Z$,
    yet $[X, Y] = Z$.
    The central direction is therefore not a degree of freedom of the BFS
    exploration but the structural obstruction to integrability of the horizontal
    distribution --- the torsion of the sub-Riemannian geometry of
    $\mathrm{Heis}_3$.
    The admissibility form $\mathcal{E}_q$, built on the perturbation of admissible
    states by $\rho_q(s)$, is directly sensitive to this non-commutativity.
    In the large-$q$ Mosco limit~\cite{Beau2026Q5a}, this torsion translates into
    the Laplace-type operator $L_\Pi = A_g$ that enters the present paper.
    The Einstein response derived here is therefore, at the structural level,
    a consequence of Heisenberg torsion rather than of a postulated spacetime
    curvature.
  \end{remark}

  Within the spectral framework considered in this work, the geometric response
  arises from the renormalized variation of the functional
  \[
    S_\Pi[g] = \frac12 \log \det{}' A_g .
  \]
  The heat-kernel analysis shows that the infrared-dominant local contribution
  to the metric variation is proportional to the Einstein tensor.
  The spectral construction therefore provides an explicit realization of the
  effective dynamics compatible with a non-injective projective structure.

  Taken together, these results suggest the following structural picture.
  Non-injective projection allows the emergence of non-trivial fiber curvature,
  which in turn enables an effective geometric response.
  The spectral entropy functional provides a minimal mechanism realizing this
  response and selects the Einstein sector as the dominant infrared contribution
  to the effective dynamics.
  This mechanism is consistent with the relational spectral framework
  of~\cite{Beau2026a}, in which spacetime geometry emerges operationally from
  relational configurations under non-injective projection.
  The same non-injectivity argument further yields gauge charges and gauge groups
  as fibre invariants of $\Pi$, establishing Yang\textendash Mills structure without
  additional postulates~\cite{Beau2026m}.
